\subsection{Perpendicular-bisector control}\label{sec:bisector_protocol}
This additional qualitative control is excluded from the five-sequence, 614-frame aggregates. For $t=0,\ldots,63$, the field is the sum of three isotropic Gaussians with standard deviation 4, amplitudes $1$, $1.05$, and $0.95$, and centers $\mathsf A=(20,20)$, $\mathsf B(t)=(40+40t/63,80-40t/63)$, and $\mathsf C=(100,100)$. We sample $[0,120]^2$ on a $241\times241$ grid and extract augmented split trees without an added ridge. Every frame has three retained extrema. References use extracted grid extrema with $\mathbf d=(1,0)$, not the analytic centers; grid spacing is 0.5. All 64 moving-feature references and all 64 achieved RG-MTM-R anchors are distinct, with no adjacent-frame repeats. Displayed trajectories are not smoothed. The two mirror endpoints have identical feature-distance matrices.

We retain the declared RG-MTM-R interval model and $\kappa=0.20$. The observed pairwise merge-level error is zero in every frame for all three methods, so this experiment does not establish the necessity of relaxation. ST-MTM uses uniform weights, $K=65$, $r=0.95$, temporal weight $1.5$, 4096 layout samples, and start timestep 0. RG-MTM-R also uses 4096 layout samples; TMTM retains the native sample permutation. Baseline map axes display internal coordinates. Only when measuring achieved positional error, each baseline receives one sequence-wide affine minimax calibration, allowing a global reflection but no framewise alignment. These errors must not be mixed with the main suite's order-only, monotone-calibrated $P$.

At the final frame, achieved positional errors are $16.67\%$, $16.67\%$, and $0.84\%$ of $L$ for TMTM, ST-MTM, and RG-MTM-R. Corresponding extrema-based stress-1 values are $0.053$, $0.053$, and $0.264$; these are not ST-MTM's SNS. The control illustrates the additional information supplied by a fixed world reference, while disclosing the larger relative-distance stress. It is not a reproduction of the Gaussian figure in the ST-MTM paper.

\paragraph{Reproducibility.}
The experiment entry point is:\\
\path{experiments/publication2027/bisector_control.py}\\
Invoke it with arguments \texttt{bisector\_highres} and \texttt{x}. Source fields, extracted positions, rendered maps, metrics, and source hashes are stored separately under \path{results/publication2027/bisector_control/}. The main figure reuses the compact rendering geometry and consistent feature-identity colours. Baseline-validation records describe a comparison of the original TMTM C++ greedy core with the local implementation on shared extracted trees; this is not a full Inviwo/TTK pipeline run.
